% Chapter 10: Observability, Drift Monitoring & Active Learning

\begin{frame}[plain]
    \vfill
    \centering
    {\footnotesize \textbf{\textcolor{deepblue}{\textsc{TaxonVision-MLOps Architecture}}}} \\[0.35cm]
    {\large \textbf{\textcolor{signalorange}{CHAPTER 10}}} \\[0.25cm]
    {\huge \textbf{Observability \& Active Learning}} \\[0.35cm]
    {\normalsize \textcolor{black!70}{Prometheus Telemetry, Concept Drift Detection \& Continuous Human-in-the-Loop Triage}} \\[0.5cm]
    \textcolor{deepblue}{\rule{0.55\linewidth}{1.2pt}} \\[0.5cm]
    \statuspartial
    \vfill
    \script{In Chapter 10, we present our observability and active learning architecture. We must be completely candid about implementation reality: while our core inference pipeline and basic metric endpoints are operational, full drift monitoring, Evidently AI, Prometheus/Grafana scrape daemons, and Label Studio triage loops are high-priority target roadmap items currently in progress or planned for upcoming sprints.}
    \qa{How does MLOps observability differ from traditional software monitoring?}{An ML API can return HTTP 200 OK while failing silently; we monitor average conformal prediction set sizes and drift distances to detect data shifts in real time.}
\end{frame}

\begin{frame}{Observability}{Production Telemetry: Prometheus Metrics \& Grafana (Roadmap)}
    \begin{block}{Domain-Specific Telemetry Beyond HTTP 200}
        \begin{itemize}\small
            \item \textbf{System Metrics:} p95/p99 inference latency ($<50$ ms budget), GPU VRAM utilization \textcolor{signalorange}{\textbf{[Target Metric]}}.
            \item \textbf{Epistemic Metrics:} Average conformal set size ($|C(X)|$), energy-based OOD anomaly rate \textcolor{signalorange}{\textbf{[Target Metric]}}.
            \item \textbf{Operational Metrics:} Triage referral rate to human taxonomists (Label Studio queue depth) \textcolor{deepblue}{\textbf{[Roadmap]}}.
            \item \textbf{Instrumentation Status:} FastAPI \texttt{/metrics} endpoint and basic counters implemented in \texttt{telemetry.py} \textcolor{plantgreen}{\textbf{[Implemented]}}; external Prometheus scraper and Grafana dashboards \textcolor{deepblue}{\textbf{[Target Roadmap]}}.
        \end{itemize}
    \end{block}

    \vspace{0.15cm}
    \statustarget

    \script{\footnotesize Observability in MLOps differs fundamentally from DevOps: an API can return HTTP 200 OK while predicting nonsense. We track conformal set sizes, referral rates, and OOD energy. Transparently, while FastAPI exposes \texttt{/metrics} and core counters in \texttt{telemetry.py}, dedicated Prometheus scrapers and Grafana dashboards are target Phase 3 roadmap infrastructure.}
    \qa{Are Prometheus and Grafana actively scraping the service?}{\footnotesize Not yet. Metric counters exist in \texttt{telemetry.py}, but Prometheus daemons and Grafana dashboards remain on our roadmap.}
    \qa{Why pull-based Prometheus scraping over push?}{\footnotesize Pull scraping prevents telemetry bottlenecks or database timeouts from blocking the inference path.}
\end{frame}

\begin{frame}{Telemetry Architecture}{Target Telemetry Topology: FastAPI $\rightarrow$ Prometheus $\rightarrow$ Grafana}
    \begin{center}
    \resizebox{0.82\textwidth}{!}{
    \begin{tikzpicture}[
        node distance=1.5cm and 2cm,
        box/.style={rectangle, draw=deepblue, thick, fill=blue!5, text width=3cm, align=center, minimum height=1.5cm, rounded corners},
        arrow/.style={->, thick, >=stealth},
        db/.style={cylinder, draw=soilbrown, thick, aspect=0.25, fill=orange!10, text width=2.5cm, align=center, minimum height=1.5cm, shape border rotate=90}
    ]
        % Nodes
        \node[box] (fastapi) {\textbf{FastAPI Pod} \\ \texttt{/metrics} endpoint \\ \textcolor{plantgreen}{\textbf{[Active Endpoint]}}};
        \node[db, right=of fastapi] (prometheus) {\textbf{Prometheus} \\ Time-Series Database \\ \textcolor{deepblue}{\textbf{[Target Daemon]}}};
        \node[box, right=of prometheus] (grafana) {\textbf{Grafana} \\ Real-time Dashboards \\ \textcolor{deepblue}{\textbf{[Target Dashboards]}}};
        \node[box, below=of fastapi, fill=green!5, draw=plantgreen] (model) {\textbf{ONNX Runtime} \\ Inference Engine \\ \textcolor{plantgreen}{\textbf{[Active Service]}}};

        % Arrows
        \draw[arrow] (model.west) |- node[near start, left, align=right, xshift=-0.1cm] {Latency \& \\ OOD stats} (fastapi.west);
        \draw[arrow, dashed] (prometheus) -- node[above] {HTTP Pull (Roadmap)} (fastapi);
        \draw[arrow, dashed] (grafana) -- node[above] {PromQL (Roadmap)} (prometheus);
    \end{tikzpicture}
    }
    \end{center}

    \vspace{0.05cm}
    \statustarget

    \script{This diagram illustrates our target telemetry architecture. In the target state, ONNX Runtime passes statistics to FastAPI, Prometheus scrapes them every 15 seconds, and Grafana renders dashboards. Today, the FastAPI /metrics plumbing is in place, while the full Prometheus and Grafana cluster deployment is our next milestone.}
    \qa{Is this telemetry pipeline currently running live in production?}{The FastAPI endpoint and ONNX runtime are active, but the Prometheus collector daemon and Grafana dashboard nodes are target architecture currently being set up.}
    \qa{Why does Prometheus pull rather than the service pushing metrics?}{Pull-based monitoring prevents the application from hanging if the metrics server goes down or becomes slow, ensuring zero impact on the critical inference path.}
\end{frame}

\begin{frame}{Drift Monitoring}{Drift Metrics: Wasserstein-1 (Prototype) vs. MMD (Target Design)}
    \begin{columns}[T]
        \begin{column}{0.48\textwidth}
            \begin{block}{Wasserstein-1 Distance (Earth Mover's)}
                \vspace{-0.2cm}
                {\footnotesize $$W_1(P, Q) = \int_{-\infty}^{\infty} |F_P(x) - F_Q(x)| \, dx$$}
                \vspace{-0.2cm}
                \footnotesize
                \textbf{Complexity:} $\mathcal{O}(N \log N)$ over 1D projected scalar norms.\\
                \textbf{Code Status:} \textcolor{signalorange}{\textbf{Offline prototype}} in \texttt{src/taxon\_vision/monitoring/drift.py}; not yet wired into live streaming inference.
            \end{block}
        \end{column}
        \begin{column}{0.48\textwidth}
            \begin{block}{Maximum Mean Discrepancy (MMD)}
                \vspace{-0.2cm}
                {\scriptsize $$\text{MMD}^2(P, Q) = \mathbb{E}[k(x, x')] - 2\mathbb{E}[k(x, y)] + \mathbb{E}[k(y, y')]$$}
                \vspace{-0.2cm}
                \footnotesize
                \textbf{Complexity:} $\mathcal{O}(N^2 \cdot D)$ using RBF Gaussian kernel in RKHS.\\
                \textbf{Code Status:} \textcolor{deepblue}{\textbf{Target Design Specification}}; not yet implemented in codebase.
            \end{block}
        \end{column}
    \end{columns}

    \vspace{0.15cm}
    \statustarget

    \script{Let's be fully transparent about our drift monitoring status. We have implemented a 1D Wasserstein scalar distance function in drift.py as an offline prototype, but it is not yet hooked into our live inference pipeline. MMD is our target algorithm for periodic high-dimensional batch audits, but it is not yet implemented in code. This represents high-priority work for our upcoming engineering sprints.}
    \qa{Are Wasserstein and MMD currently running in production?}{No. Wasserstein-1 is implemented as an offline evaluation function in drift.py, while MMD and live stream drift alerts are target roadmap specifications.}
    \qa{What remains to be done for drift monitoring?}{We need to implement the MMD kernel matrix math, connect Wasserstein drift to live prediction batches, and configure alerting thresholds.}
\end{frame}

\begin{frame}{Telemetry Architecture}{Target Metric Inventory (Prometheus Roadmap)}
    \begin{block}{Production Metric Scrape Architecture}
        \footnotesize
        Pull-based collection scraping FastAPI \texttt{/metrics} endpoint every 15 seconds.
        \begin{enumerate}
            \item \texttt{taxon\_predictions\_total} \textcolor{plantgreen}{\textbf{[Implemented in \texttt{telemetry.py}]}}
            \item \texttt{taxon\_inference\_latency\_seconds} \textcolor{deepblue}{\textbf{[Target Histogram]}}
            \item \texttt{taxon\_conformal\_set\_size} \textcolor{deepblue}{\textbf{[Target Gauge]}}
            \item \texttt{taxon\_ood\_energy} \textcolor{deepblue}{\textbf{[Target Distribution]}}
            \item \texttt{taxon\_human\_referrals\_active} \textcolor{deepblue}{\textbf{[Roadmap / Target]}}
        \end{enumerate}
    \end{block}

    \vspace{0.15cm}
    \statustarget

    \script{This slide lists our target metric inventory. Currently in telemetry.py, only the prediction counter is wired up to FastAPI routes. Epistemic metrics like conformal set size and OOD free energy gauges, as well as human referral queue tracking, are defined as constants or design targets awaiting end-to-end integration.}
    \qa{Which metrics are actively collected today?}{Only prediction counts grouped by status in predict.py. Conformal set size, latency histograms, and referral rate metrics are defined as constants but await end-to-end wiring.}
    \qa{Why measure OOD energy instead of standard softmax probabilities?}{Softmax is notoriously overconfident, pushing unknown inputs to high probability classes. Helmholtz free energy provides a more reliable, unnormalized measure of how far an input is from the training manifold.}
\end{frame}

\begin{frame}{Drift Monitoring}{Target Multi-Scale Drift Strategy \& Evidently AI (Roadmap)}
    \begin{block}{Multi-Scale Distributional Drift Architecture (Roadmap)}
        \begin{itemize}\small
            \item \textbf{Environmental Shift:} Phenological seasonality and new camera sensors cause feature representations to drift from baseline distributions.
            \item \textbf{Online Marginal Tests (1D):} 1D projection with Wasserstein-1 distance for real-time alerting \textcolor{signalorange}{\textbf{[Prototype in \texttt{drift.py}}]}.
            \item \textbf{Offline Manifold Tests (Multivariate):} Maximum Mean Discrepancy (MMD) with RBF kernels evaluates complete structural manifold shifts \textcolor{deepblue}{\textbf{[Not Yet Implemented]}}.
            \item \textbf{Evidently AI Integration:} Target framework for automated scheduled visual drift reporting \textcolor{deepblue}{\textbf{[Not Yet Setup/Configured]}}.
            \item \textbf{Drift Retraining Triggers:} Sustained divergence triggering active learning extraction \textcolor{deepblue}{\textbf{[Phase 3 Target]}}.
        \end{itemize}
    \end{block}

    \vspace{0.15cm}
    \statustarget

    \script{Machine learning models degrade over time because the real world changes—we call this data drift. We want to be completely open with our teammates about where this system stands: Evidently AI is not yet installed or set up, MMD is not yet coded, and automated drift-triggered retraining is an architectural design on our Phase 3 roadmap. There is significant, impactful engineering work left to build here.}
    \qa{Is Evidently AI currently operational in TaxonVision?}{No. Evidently AI is our planned tool for drift reports and visual dashboards, but it has not been configured or integrated into our repository yet.}
    \qa{Why not run multivariate MMD on every single incoming inference request?}{Computing full-manifold MMD across high-dimensional embeddings has quadratic complexity $O(N^2)$ in sample size. We perform 1D Wasserstein distance on projected features in real time and reserve full MMD for periodic offline batch audits.}
\end{frame}

\begin{frame}{Active Learning}{BADGE Sampling \& Human Review Queue}
    \begin{block}{The Annotation Bottleneck}
        Expert taxonomy time is expensive. Randomly sampling from the unlabelled stream yields redundant, low-diversity annotations.
        \begin{itemize}
            \item \textbf{BADGE Algorithm:} Batch Active learning by Diverse Gradient Embeddings.
            \item \textbf{Dual Criteria:} It computes hypothetical gradient embeddings representing both \textit{epistemic uncertainty} and \textit{feature diversity}.
            \item \textbf{Selection:} Uses $k$-MEANS++ clustering in gradient space to select the optimal batch.
        \end{itemize}
    \end{block}

    \begin{block}{Engineering Rationale}
        By routing only the most mathematically informative samples to Label Studio, we minimise annotation cost while maximising model accuracy improvements per retraining cycle.
    \end{block}

    \vspace{0.2cm}
    \statuspartial

    \script{Annotation by expert taxonomists is expensive and slow. If you send them 10,000 photos of the same common dandelion, you waste their time. BADGE solves this by analyzing the gradients that each unlabelled image would induce in the model. It uses k-MEANS++ clustering in gradient space to pick a batch of images that are both highly uncertain AND visually diverse. This maximizes the model's accuracy improvement per dollar spent on human labeling.}
    \qa{What is the current status of the active learning pipeline?}{The BADGE mathematical sampling algorithm is fully implemented in \texttt{src/taxon\_vision/active\_learning/badge.py}. However, the human review queue integration and automated retraining triggers are in progress.}
\end{frame}

\begin{frame}{Active Learning}{Target Active Learning Loop: Label Studio \& DagsHub (Roadmap)}
    \begin{block}{Architectural \& Operational Target Blueprint}
        \begin{itemize}\small
            \item \textbf{Current Implementation Reality:} Label Studio is \textbf{not yet set up or prepared} in our infrastructure. This slide details the target design for our Phase 3 human-in-the-loop triage engine.
            \item \textbf{Triggering Mechanisms:}
            \begin{itemize}\footnotesize
                \item \textbf{Conformal Ambiguity Filter:} Samples where prediction set size $|\hat{C}(x)| > 1$ or empty sets $|\hat{C}(x)| = 0$.
                \item \textbf{BADGE Diversity Sampling:} Clustering hallucinated gradient embeddings to capture boundary uncertainty \textcolor{plantgreen}{\textbf{[Algorithm Implemented]}}.
            \end{itemize}
            \item \textbf{Annotation Loop (Target):} Automated batch export to Label Studio on DagsHub $\to$ Domain expert review $\to$ DVC dataset commit $\to$ Jenkins retraining webhook \textcolor{deepblue}{\textbf{[Roadmap]}}.
            \item \textbf{Cost Efficiency:} Prevents sending redundant high-confidence common taxa to human reviewers, cutting labeling expenses by $> 70\%$.
        \end{itemize}
    \end{block}

    \vspace{0.15cm}
    \statustarget

    \script{\footnotesize On this slide, we illustrate our target active learning loop. To be completely clear with the team: Label Studio is not yet set up, configured, or integrated into our deployment. The sampling math in BADGE exists in badge.py, but the web UI, annotation ingestion pipelines, and automated DVC retraining webhooks are target roadmap items that our team will build next.}
    \qa{Is Label Studio currently deployed and operational?}{\footnotesize No. There is nothing prepared for Label Studio yet in our infrastructure; it is purely a target architectural blueprint for Phase 3.}
    \qa{Why show this workflow if Label Studio is not yet set up?}{\footnotesize To define the exact data contract and trigger mechanisms so that the engineering team can build the Label Studio connector against a unified specification.}
\end{frame}

\begin{frame}{Active Learning Flow}{Target Closed-Loop Architecture: Stream $\rightarrow$ BADGE $\rightarrow$ Label Studio}
    \begin{center}
    \resizebox{!}{0.48\textheight}{
    \begin{tikzpicture}[
        node distance=0.8cm and 1.2cm,
        stepbox/.style={rectangle, draw=deepblue, thick, fill=blue!5, text width=2.5cm, align=center, minimum height=1.0cm, rounded corners},
        arrow/.style={->, thick, >=stealth},
        labelbox/.style={rectangle, draw=signalorange, thick, fill=orange!10, text width=2.5cm, align=center, minimum height=1.0cm, rounded corners}
    ]
        % Nodes
        \node[stepbox] (stream) {\textbf{1. Live Stream} \\ Unlabelled Data \\ {\footnotesize\textcolor{plantgreen}{\textbf{[Active Data]}}}};
        \node[stepbox, right=of stream] (conformal) {\textbf{2. Conformal Filter} \\ Identify ambiguous \\ {\footnotesize\textcolor{plantgreen}{\textbf{[Active Filter]}}}};
        \node[stepbox, right=of conformal] (badge) {\textbf{3. BADGE} \\ Diversity sampling \\ {\footnotesize\textcolor{plantgreen}{\textbf{[Implemented]}}}};
        \node[labelbox, below=of badge] (labelstudio) {\textbf{4. Label Studio} \\ Expert Review \\ {\footnotesize\textcolor{deepblue}{\textbf{[Target System]}}}};
        \node[stepbox, left=of labelstudio] (retrain) {\textbf{5. Retraining} \\ DVC Retrain Trigger \\ {\footnotesize\textcolor{deepblue}{\textbf{[Target Webhook]}}}};

        % Arrows
        \draw[arrow] (stream) -- (conformal);
        \draw[arrow] (conformal) -- (badge);
        \draw[arrow, dashed] (badge) -- (labelstudio);
        \draw[arrow, dashed] (labelstudio) -- (retrain);
        \draw[arrow, dashed] (retrain) -| (stream);
    \end{tikzpicture}
    }
    \end{center}

    \vspace{0.05cm}
    \statustarget

    \script{\footnotesize This flow shows our target architecture for active learning in production. While streaming observations and conformal filtering are active, and BADGE sampling math is implemented in badge.py, automated Label Studio dispatch and closed-loop retraining webhooks remain roadmap goals.}
    \qa{Which parts of this 5-step loop are implemented today?}{\footnotesize Steps 1 and 2 (streaming data and conformal filtering) and step 3 math (BADGE) are ready. Steps 4 and 5 (Label Studio annotation and automated retraining webhook) are roadmap items.}
    \qa{Why sequence Conformal Filtering before BADGE?}{\footnotesize Conformal filtering quickly isolates ambiguous predictions, shrinking the candidate pool so that BADGE gradient clustering runs in seconds.}
\end{frame}
